\begin{abstract}
\noindent\textbf{Motivation:} Clustering is a foundational step in single-cell RNA
sequencing (scRNA-seq) analysis, supporting both the identification of common cell
types and the discovery of rare populations. Deep embedded clustering (DEC) couples
the clustering objective with representation learning, yet in comprehensive benchmarks
DEC-style models rank near the bottom, and the reason for this failure has not been
characterized.

\noindent\textbf{Results:} We identify a concrete failure mode of DEC on scRNA-seq
data: per-dimension variance collapse of the clustering latent. The confidence-gated
KL sharpening that drives DEC compresses the embedding onto a small number of active
coordinates, reducing the effective dimensionality from 114 to 60 and degrading
$k$-means read-out. Building on a masked-autoencoder (scMAE) backbone with zero-masking,
we add a single per-dimension variance-floor term, borrowed from VICReg, that penalizes
any latent coordinate whose batch standard deviation falls below one. This term restores
the effective dimensionality to 114, improves the Adjusted Rand Index (ARI) over plain
DEC by $0.272$ on average across 18 datasets (winning on 17), improves Normalized Mutual
Information by $0.238$ (winning on all 18), and lowers seed-to-seed variance by a factor
of $2.5$. We state explicitly that the resulting model, VarFloor-scMAE, matches rather
than surpasses the accuracy of the scMAE backbone; its contribution is diagnostic and
corrective with respect to DEC. A re-evaluation under macro-F1 and rare-cluster recall
shows that ARI systematically over-credits sharpening methods, and that rare-cell
discovery remains an open problem that no method examined here, including ours, resolves.

\noindent\textbf{Availability:} Code and processed data are available on request.
\end{abstract}
